\section{Cyclic lifts with the same classes}\label{sec:constructions}

The original blow-up construction in \cite{ClassUniformMemo2026} gives
unbounded absolute excess over the interval bound with a growing number
of classes. Its revised memo also gives a two-class parity-coset family
for every $d=4k\ge8$ over fields where $2\ne0$
\cite[Theorem~4.8]{ClassUniformMemo2026}. The revised seventeenth paper
extends that rule to every field, including characteristic two, and uses
mixed blow-ups to obtain unbounded excess for every fixed number of
classes\cite{SeventeenthPaper2026}.
The construction below instead gives a lift for an arbitrary seed with
the same number of classes, an exact periodic-capacity formula, and a
paid realization of its exhibited binary rule. These fixed-class families
all have vanishing normalized excess; fixed-class or binary additive
excess alone is not a new conclusion here.

\begin{theorem}[Deterministic cyclic lift]\label{thm:cyclic-lift}
\formal{ConstructionLinearLift.lean}{cyclic_lift}
Let a linear source $C_s\le(F^b)^{s\times s}$ have dimension $D$ and
deterministic shared decoders for an $M$-class partition $p$ of $\Z/s$.
Let $m\ge1$ be coprime to $s$, with $m\ne0$ in $F$, and put $d=ms$.
The lifted partition $p^\uparrow(c)=p(c\bmod s)$ admits a linear source
of dimension
\begin{equation}\label{eq:cyclic-lift-dimension}
 D^\uparrow=\frac{b(d^2-ms^2)}2+mD
\end{equation}
with deterministic shared decoders. Linear seed decoders give linear
lifted decoders.
\end{theorem}
\begin{outline}
Use Chinese-remainder coordinates $(a,r)\in\Z/m\times\Z/s$ for both
rows and columns. For $z\in\Z/m$ define
\[
 Z_z(X)_{r,c}=\frac1m\sum_{a\in\Z/m}X_{(a,r),(z-a,c)},\qquad
 (EX)_{(a,r),(h,c)}=Z_{a+h}(X)_{r,c}.
\]
The map $E$ is an idempotent. Its image consists of arrays whose outer
indices enter only through their sum, and is parametrized by the $m$
independent seed matrices $Z_z$. Thus $\rank E=mbs^2$.
It commutes with column rotation and transpose. On its image transpose
acts as the inner $s\times s$ transpose on each $Z_z$.

Take the direct sum
\[
 C^\uparrow=(\ker E\cap\operatorname{Fix}\tau)\oplus
 \{X\in\operatorname{im}E:Z_z(X)\in C_s\text{ for every }z\}.
\]
The second summand has dimension $mD$. The symmetric space has dimension
$bd(d+1)/2$ in every characteristic. Its image under $E$ is exactly the
space of $m$ symmetric seed matrices, of dimension $mbs(s+1)/2$:
commutation proves inclusion, and every such image vector is already
symmetric and fixed by $E$. Rank-nullity gives the first summand's
dimension $b(d^2-ms^2)/2$.

For an arbitrary full view $Y$, define one decoder per original class by
\begin{equation}\label{eq:lifted-deterministic-decoder}
 F_J(Y)_{a,r}=(Y-EY)_{(a,r),(0,0)}+f_J(Z_a(Y))_r.
\end{equation}
At a receiver with CRT index $(u,v)$, literal rotation satisfies
$Z_a(T^{(u,v)}X)=T_s^vZ_{a+u}(X)$. Thus the complete induced map is
\[
 K^\uparrow(X)_{(a,r),(u,v)}
 =(X-EX)_{(a,r),(u,v)}+f_{p(v)}(T_s^vZ_{a+u}(X))_r.
\]
For a codeword, the first term is its symmetric part, and the second
transposes its projected part. This proves correctness. The receiver
index is not an extra argument to \eqref{eq:lifted-deterministic-decoder};
only its seed class selects the function. In particular a nonlinear
decoder is applied after averaging the seed matrix, never distributed
through that average. All displayed operations are linear when the
seed decoder is linear.
\end{outline}

For linear rules there is also an operator description. On the outer
tensor square put
$E=m^{-1}\sum_{t\in\Z/m}G^t\otimes G^{-t}$, where $G$ is rotation.
After regrouping outer and inner factors, the complete lifted rule is
\begin{equation}\label{eq:lift-operator}
 K^\uparrow=(I-E)\otimes I+E\otimes K_s.
\end{equation}
If a seed term has column factor $P_JS^e$, its lifted column factor is
$G^{-t}\otimes P_JS^e=P_{J^\uparrow}\sigma_d^u$, where CRT gives
$u\equiv-t\pmod m$ and $u\equiv e\pmod s$. Hence this is one allowed
same-class operator, not a composition with a free coordinate permutation.
If the seed code is the entire agreement kernel, the displayed lifted
code is the entire kernel of \eqref{eq:lift-operator}.

\begin{corollary}[A fixed two-class excess]\label{cor:fixed-two-excess}
\formal{ConstructionSeed.lean}{binaryCyclicCode_linear, binaryCyclicPartition_surjective, finrank_binaryCyclicCode, binaryCyclicCode_interval_excess}
For every $m\ge1$ coprime to six, $d=6m$ admits a binary two-class
linear rule and a linear source of dimension
\[
 \frac{d^2}{2}+\frac{7d}{6}.
\]
Its excess over the two-interval bound is $d/6$.
\end{corollary}
\begin{outline}
Use the published six-dimensional seed with classes
$\{0,1,3\},\{2,4,5\}$ and its verified $25$-dimensional code
from \cite{SeventeenthPaper2026,ClassUniformMemo2026}.
Denote the complete rule induced by these seed decoders by $K_6$.
Since $m$ is odd, it is invertible in $\F_2$.
Equation \eqref{eq:cyclic-lift-dimension} gives
$18m^2+7m=d^2/2+7d/6$; subtract $d^2/2+d$.
The conclusion needs a seed subcode of dimension $25$, not a proof that
$25$ is the seed's maximum capacity.
\end{outline}

\subsection{The exact linear capacity of a periodic lift}

Write $\operatorname{Cap}^{\mathrm{lin}}_F(p)$ for capacity with linear
decoders over $F$.

\begin{theorem}[Periodic capacity reduction]\label{thm:periodic-capacity}
\formal{ConstructionCapacity.lean}{periodic_linear_capacity_split, periodic_linear_capacity_lower, periodic_linear_capacity_algebraicClosure_upper}
In the scalar case $b=1$, suppose additionally that $F$ contains all
$m$-th roots of unity. Then
\[
 \operatorname{Cap}^{\mathrm{lin}}_F(p^\uparrow)
 =\frac{d^2-ms^2}{2}+m\operatorname{Cap}^{\mathrm{lin}}_F(p).
\]
Without the splitting assumption, the same expression with $F$ is a lower
bound, and the expression with its algebraic closure $\overline F$ is an
upper bound.
\end{theorem}
\begin{outline}
The lower bound is Theorem \ref{thm:cyclic-lift}. For the upper bound,
the allowed column space in CRT coordinates is
$\operatorname{Circ}_m\otimes\operatorname{span}\{P_JS^e\}$.
Simultaneously diagonalize the outer circulants, using the same basis
on both tensor factors so that the demand remains swap. This is valid
because $m$ is invertible and its roots split. Write the input as an
$m\times m$ array of $s\times s$ blocks $X_{ij}$. Correctness has the form
\[
 X_{ji}^{\mathsf T}=\sum_a L_{ia}^j(X_{aj}),
\]
where each $L_{ia}^j$ is an allowed seed rule. In particular
$\delta_j=\dim\ker(\tau_s-L_{jj}^j)\le
\operatorname{Cap}^{\mathrm{lin}}_F(p)$.

Supply all upper-triangular blocks $X_{ij}$, $i<j$, and process $j$ in
decreasing order. All off-diagonal blocks of column $j$ are known: those
above were supplied, those below belong to recovered rows. The equation
with $i=j$ determines $X_{jj}$ up to a vector in a fixed $\delta_j$-
dimensional kernel. Supply $\delta_j$ linear coordinates injective on
that kernel; this determines $X_{jj}$ uniquely even on an affine fibre.
The other equations now recover row $j$. This is an injective linear
description of size $s^2m(m-1)/2+\sum_j\delta_j$, proving the bound.
For nonsplitting $F$, embed a basis of a maximizing source into
$\overline F$ and take its $\overline F$-span. Scalar extension preserves
the basis's independence and extends every class rule by its coordinate
coefficients. The extended source has the same dimension and the same
receiver classes. Apply the split result over $\overline F$.
\end{outline}

Thus every rule on a fixed periodic seed partition has rate
$1/2+O(1/d)$, even if it is not the exhibited lift. The qualitative limit
already follows from the period obstruction in \cite{SeventeenthPaper2026};
the equality records the exact finite seed-capacity contribution.

\subsection{A paid implementation of the exhibited binary lift}

\begin{theorem}[Two-tape implementation]\label{thm:paid-lift}
\notformalized
There is a fixed finite-alphabet, finite-control machine with two initially
blank work tapes, read-only input and append-only unreadable output that
computes the binary operator in Corollary \ref{cor:fixed-two-excess}
lane by lane in $O(P)$ time, $P=d^2b$, on admitted dimensions $d=6m$,
$\gcd(m,6)=1$. There are separate machines for $b=1$ and for
$b=\max(1,\lceil\log_2d\rceil)$. Each has a transpose-correct subcode of
dimension $b(d^2/2+7d/6)$.
\end{theorem}
\begin{informalproof}
Here is a complete bounded-state schedule, including preparation. Put
$\mu=mb$. The first work tape has $m$ consecutive $b$-cell words and
$36$ binary payload tracks per physical cell. For $z\in\Z/m$ their
intended contents are
\[
 H_z[i,j]=\bigoplus_{r\equiv i\ (6),\ c\equiv j\ (6),\ r+c\equiv z\ (m)}X_{r,c},
 \qquad 0\le i,j<6.
\]
Each word has independent FIRST/LAST marker bits; HOME is at cell zero,
END at $\mu$, and START at one word's first cell. Coincident markers
when $b=1$ or $m=1$ are distinct alphabet tracks. All payload starts zero.
After preparation the second work head parks throughout both data passes.

At the start of input row $r$, START and the bank head are at word
$r\bmod m$, with finite registers $i=r\bmod6$, $j=0$, and lap count zero.
A DATA step reads the old raw bit and old bank symbol, advances both heads
right, and either toggles $H[i,j]$ by that bit (pass one) or appends
$x+H[i,j]+(K_6H)[i,j]$ (pass two), preserving the bank. The last expression
is a fixed Boolean function of $36$ old bits. An OLD LAST marker increments
$j$ modulo six and enters LAND. LAND at END rewinds physically to HOME;
else it is at a FIRST marker. Returning to START counts a completed lap.
Six laps finish one row. Importantly $j$ is not reset at a lap boundary:
$m$ need not be divisible by six. At row end clear START, physically advance
one whole word, wrapping at END if necessary, and set START there; then
increment $i$ modulo six and reset $j$ and the lap counter. Perform this
offset after the last row too. Only the row-start mode tests the raw
endmarker; housekeeping preserves payload, holds raw still, and emits
nothing.

Induction gives bank address $((r+c)\bmod m)b+t$ at lane $t$ of entry
$(r,c)$, with $j=c\bmod6$. Thus pass one computes exactly $H$. After
$d$ rows its START is HOME. Rewind raw across its entire input, preserving
the bank, and run the same schedule in emission mode. There are exactly
$P$ appends in ordinary row-major order. Because $m$ is odd, the sum in
$H$ is precisely the binary averaging projector $E$; the emitted formula
is \eqref{eq:lift-operator}. All reads are pretransition reads, and all
writes are at the old head locations.

For completeness one pass has $P$ DATA steps, $d^2$ ordinary LAND steps,
and six physical END wraps per row, costing $6d(\mu+1)$ extra steps.
Row offsets cost $db+d+6(\mu+1)$, and row-start branches cost $d$.
The total is
\[
 2P+d^2+2db+8d+6\le5P\qquad(d\ge6,b\ge1).
\]
The between-pass raw rewind costs $P+2$; the final absent transition halts.

It remains to establish this bank from blanks without size advice. Count
the raw length on the second tape using a little-endian binary counter
and a marked least-significant cell. A carry through $k$ ones and return
costs $O(k+1)$, whose total over $P$ increments is $O(P)$; hold raw fixed
during each carry, then advance it. Rewind raw. Fixed schoolbook arithmetic
on a constant number of delimited $O(\log(P+2))$-bit registers verifies
$P=d^2$ or $P=d^2\max(1,\lceil\log_2d\rceil)$ by binary search, and
checks $d=6m$, $\gcd(m,6)=1$. These functions are strictly increasing in
$d$, and all arithmetic takes $\operatorname{polylog}(P+2)$ time.
Reject other lengths before output.

The first head now writes the bank sequentially while the second head
maintains counters for $m$ words and $b$ cells. Binary decrement and return
has aggregate $O(b+\log(P+2))$ cost per word; no full-width zero test is
performed per cell. Setup, copying and register trips cost at most
$\operatorname{polylog}(P+2)$ per word. Install all markers, write END,
return to HOME, and erase the second tape's delimited scratch frame.
This costs
$O(P+m\log^c(P+2)+mb+\log^c(P+2))=O(P)$ for a fixed exponent $c$,
since $m=d/6$ and $d\log^c(d+2)=O(d^2b)$ in both width regimes.
Every unbounded counter was tape data, and all other registers range over
fixed finite sets. This establishes the promised initial boundary with
exactly two work heads and proves the uniform time bound.
\end{informalproof}

This is an implementation of one displayed operator on all inputs of its
admitted lengths. Its transpose-correct source has rate tending to one
half, so it is not a full-input transposer or a lower bound against native
machines. No general class-uniform operator is assigned a free implementation.
